\section{因果逐时求解、滚动状态与动态复核}
\label{sec:solution}

\subsection{MILP求解入口}

在给定当前设备边界、初始状态和管理目标后，联合问题写为
\begin{equation}
\min_{\bm x_t,\bm y_t} F_{eco,t}
\quad\mathrm{s.t.}\quad
\bm A_t\bm x_t+\bm B_t\bm y_t\le\bm b_t,
\quad
\bm A_t^{eq}\bm x_t+\bm B_t^{eq}\bm y_t=\bm b_t^{eq},
\label{eq:milp-compact}
\end{equation}
其中$\bm x_t$为连续变量，$\bm y_t$为整数/二元变量；约束矩阵由4.9的物理、逻辑、
服务和目标底线组装。若主目标为可靠性或消纳率，则采用词典序求解：先固定第一层最优值，
再在不恶化该值的可行域中优化下一层，避免任意加权系数改变服务优先级。

\subsection{因果信息边界}

实际应用不得把未来真实序列作为已知量。决策必须满足非前视条件
\begin{equation}
\bm u_t=\pi_\theta(\mathcal I_t),\qquad
\mathcal I_t=\sigma\!\left(
\{\bm w_\tau,\bm d_\tau,\bm p_\tau,z_\tau^{evt}\}_{\tau\le t},
\bm s_{t-1},\Theta^{contract}\right),
\label{eq:causal-policy}
\end{equation}
其中$\bm w_t$为当前源侧观测，$\bm d_t$为当前负荷/任务，$\bm p_t$为当前价格，
$\bm s_{t-1}$为上一小时末电池、氢库存和设备状态，$\Theta^{contract}$为事先已知的合同
与设备参数。$\mathcal I_t$不包含$t+1$之后的真实源、真实价格、真实负荷或全年总供电量。

现行在线策略把计划信息集与结算观测分开。计划动作满足
\begin{equation}
\bm u_t^{plan}=\pi_\theta(\mathcal I_t^-),\qquad
\mathcal I_t^-=\sigma\!\left(
\{\bm w_\tau\}_{\tau\le t-1},z_t^{evt},\bm d_t,\bm p_t,
\bm s_{t-1},\Theta^{contract}\right),
\label{eq:causal-planning-policy}
\end{equation}
计划侧不读取$\bm w_t$。结算时读取当前实测$\bm w_t$，在计划份额容差和物理约束下
形成$\bm u_t$；结算完成后才用$\bm w_t$更新预测器。卡尔曼递推采用经典线性滤波
结构\cite{kalman1960}；小时类型先验与动态后验来自重叠历史，使用协方差交集处理未知
交叉相关\cite{julier1997ci}。现行执行链为
\begin{enumerate}[label=(\arabic*)]
\item 用历史多源实测和当前事件状态形成当前小时源侧先验--后验预测；
\item 继承$t-1$末BESS能量、氢库存、电解槽/算力/氢转电状态；
\item 在预测功率下求原始E/H/C计划，并用平滑、死区与小时变化限幅抑制抖动；
\item 读取当前真实多源功率，先尝试同时满足零ENS、计划份额和BESS储备目标；失败时依次释放产品份额与储备约束，无储备约束仍不可行才按$\varepsilon$约束法先最小化ENS、再在最小ENS容差内优化经济目标\cite{mavrotas2009}；
\item 仅当求解状态和独立审计均通过时提交动作，并将时段末状态传至$t+1$。
\end{enumerate}
该方法避免完美预见，并以可回退SOC储备约束缓解单小时视野的短视库存消耗；但仍不能替代多时段预测。后续可扩展为6、24或48~h
滚动MPC：预测必须由当时可获得的历史、数值预报或合同计划生成，每次仍只执行首时段，
且不得用事后真实数据替代预测。

\subsection{资产基准、在线策略与离线基准的关系}

当前比较集合只有纯电、纯氢、纯算三类资产基准和一个在线先验--后验策略。资产基准通过
关闭另外两类可选产品资产形成，不强制E/H/C比例；原22种固定比例方案已经退出比较与
代码。四方案按统一资源、设备、价格、负荷、初始状态和事件时序报告经济、消纳、可靠性
和状态轨迹。主比较把海底算力全部设为可中断柔性负荷；不可中断SLA只在另行敏感性中启用。

若另行求解已知完整48~h或年度真实序列的全时域MILP，该结果只能作为离线“完美信息上界”
或算法诊断基准，必须明确标注，不得当作可部署策略。正式运行结论以满足
式~\eqref{eq:causal-policy}的信息边界为准。

\subsection{极端事件策略}

台风等事件采用“预警--过境--恢复”状态机。预警在当前时段可知时，暂停可选海缆外送、
制氢和柔性算力，利用可用源功率恢复BESS；过境期关闭不具备生存或通信
条件的设备，优先保障内部关键负荷和刚性服务，并在明确启用的韧性配置中允许4.5氢转电；
恢复期按实测可用率和启动限制逐步复位。BESS正常、预警和恢复储备目标均以可回退的
单时段末端约束实现，不能优先于当前零ENS。氢库存只保留物理非负与容量上限，不设置运营
储备下限；事件前不得在模型外增加氢或SOC。事件规则只代表调度应对，不能证明漂浮平台、海缆或海底数据舱具备台风结构
生存性。

\subsection{求解审计与多时间尺度复核}

结果可接受需同时满足：求解器返回可用整数解；全部输入和输出有限；式
~\eqref{eq:bus-residual}、电池能量和储氢质量残差在容差内；容量、SOC、库存、互斥、
逐时比例和服务约束均通过；事件策略和资产开关未被绕过。求解失败、超时无可用解或任一
硬约束审计失败时，该时段不提交，不能以启发式结果冒充最优解。

小时调度结果还须经过以下动态门：
\begin{enumerate}[label=(\arabic*)]
\item 5--15~min功率可行性复核源侧快速爬坡、PCS和电解槽动态；
\item REGFM正序模型检查构网有功/无功、限流与频率响应\cite{regfm-b1}；
\item EMT与HIL检查电压、故障穿越、保护配合、谐波和多机并联；
\item 黑启动、船运窗口、海工生存和现场试验完成工程验收。
\end{enumerate}
动态复核不通过时，应收紧SOC、备用、P--Q、爬坡或端口容量后返回小时模型重算，而不是
在报告中删除不利工况。本节只给出求解与复核流程，不列示48~h或年度测试结果。
